\CUSTOMstruct{ЗАКЛЮЧЕНИЕ}

В ходе учебной практики выполнена подготовка к написанию выпускной квалификационной работы: отработаны навыки работы с проектной и технической документацией, сформирован отчёт о выполненных этапах.

Прослушаны два лекционных занятия. На первом зафиксированы виды и структура проектной документации, требования к ВКР и техническому заданию (в том числе по документу ЛНАОБУЧ-СМК-03-05-2022), 
стандарты документирования и типовую структуру ТЗ. На втором занятии получена обратная связь по отчёту: учтены замечания преподавателя по введению (таблица этапов), изложению от первого лица, 
обоснованию выбора LaTeX и описанию сделанной работы с приложением итогов.

Разработан шаблон для ВКР в LaTeX с опорой на разделы 4.3--4.11 требований к оформлению: реализованы нумерация страниц, заголовки, иллюстрации, таблицы, формулы, приложения и переиспользуемые команды оформления. 
Шаблон позволяет собирать документы, соответствующие единым правилам университета.

Составлено техническое задание на систему «Разработка гибридной системы интеллектуального нормоконтроля с интеграцией LLM для семантического анализа технической документации». 
При наполнении ТЗ использованы научные источники по согласованию документов, BPMN и автоматизации бизнес-процессов; функциональные требования выведены из use case диаграммы, 
назначение и актуальность --- из опроса, структура системы --- из схемы в Eraiser, стадии --- из архитектуры и этапов защиты ВКР. Документ оформлен в LaTeX с переиспользованием шаблона и приведён в приложении~Б.

В результате подготовлены основа для оформления ВКР и проектная документация, что соответствует целям практики по этапам 1.1--1.3.

